The accurate identification and monitoring of extreme weather events, such as tropical cyclones (TCs) and atmospheric rivers (ARs), are vital for modern meteorology  \cite{hewson2010objective, hengstebeck2018radar, dawe2012statistical, lawrence2018characterizing, rautenhaus2017visualization, beckert2023three}. These phenomena significantly impact global weather patterns, and their precise detection is a cornerstone for improving forecasting accuracy and understanding climate dynamics. While traditional detection methods have relied on rule-based systems using physical heuristics—such as sea-level pressure minima for TCs or specific moisture thresholds for ARs—these approaches often struggle with consistency and universal formulation  \cite{neu2013imilast, bourdin2022intercomparison, guan2015detection, shields2018atmospheric}.

Recent advancements in deep learning, particularly convolutional neural networks like U-Net and CGNet, have transformed these detection tasks into semantic segmentation problems. These models have demonstrated an ability to learn physically plausible, complex atmospheric patterns that are difficult to define through rigid mathematical rules alone \cite{boukabara2021outlook, higgins2023using}. However, as the field moves toward more robust and versatile architectures, the use of foundation models to further enhance detection capabilities still have significant research gaps.

\subsection{Adapting SAM to the Climate Domain}

The Segment Anything Model (SAM) has emerged as a powerful tool in computer vision due to its remarkable zero-shot generalization. Despite its success with natural images, applying SAM to meteorological data presents several unique challenges:

\begin{itemize}
\item Data Dimensionality and Representation: SAM is natively designed for three-channel RGB inputs. In contrast, the ClimateNet dataset consists of 16-channel atmospheric grids, including variables like zonal wind, surface temperature, and precipitable water. These variables lack the high-frequency edges and clear semantic boundaries found in natural imagery, necessitating specialized adaptation of the image encoder.
\item Transition from Binary to Multiclass Segmentation: By design, SAM is a prompt-driven binary segmentation model that distinguishes a foreground object from the background based on a user-provided prompt. To make it suitable for climate science, the architecture must be modified to handle multiclass semantic segmentation, enabling the simultaneous and distinct detection of TCs and ARs within a single global raster.
\item Autonomous Prompt Generation: SAM's reliance on manual geometric prompts is a significant bottleneck for real-time weather forecasting. To make the system operationally viable, it must transition toward a prompt-free architecture. This requires the development of internal mechanisms—such as multi-scale fusion or auxiliary prediction networks—to automatically generate the guidance previously provided by humans.
\end{itemize}

\subsection{Aim of the Thesis}

The primary objective of this thesis is to develop ClimateSAM, a framework that adapts the Segment Anything Model for the automated, multiclass detection of atmospheric features in the ClimateNet dataset. This research aims to bridge the gap between foundation models and climate science by investigating parameter-efficient fine-tuning (PEFT) techniques, such as Low-Rank Adaptation (LoRA) and learnable token-based adapters, to refine SAM's encoder for non-RGB atmospheric data. Furthermore, the study explores the integration of automated prompt generators to eliminate the need for manual intervention, thereby creating a scalable pipeline for identifying extreme weather events. Ultimately, this work seeks to evaluate whether such an adapted foundation model can capture the unique morphological structures of tropical cyclones and atmospheric rivers more effectively than existing baselines.